%% kapa: exam=1 round=30 year=2019 period=1 form=A subject=economics subject_ko=경제학원론 seq=3 number=43
%% source: Q-Net 감정평가사 기출문제 (cst00309) artlSeq=5205205 file=extracted/1008/5205205/2019년 제30회 감정평가사 1차시험 문제지/2019년 제30회 감정평가사 1차시험 문제지/2019년도 제30회 감정평가사 1차 1교시 A형.hwp
%% conversion: kordoc 4.24.51 -> tools/gen1.py (auto); review=proofread-sonnet; flags=none
\begin{question}{43}
소비자 甲은 ${ X }$재와 ${ Y }$재만 소비하여 효용을 극대화한다. 제1기의 ${ X }$재 가격은 3이고, ${ Y }$재 가격은 6이었을 때, 소비조합 (${ X=3, Y=5 }$)를 선택하였다. 제2기에는 동일한 소득에서 ${ X }$재와 ${ Y }$재의 변동된 가격 ${ P_{ X } }$, ${ P_{ Y } }$에서 소비조합 (${ X=6, Y=3 }$)을 선택하였다. 甲의 선택이 현시선호 약공리(weak axiom)를 만족하기 위한 조건은?
\begin{choices}
  \item $2P_{ X } < 3P_{ Y }$
  \item $2P_{ X } >3P_{ Y }$
  \item $3P_{ X } <2P_{ Y }$
  \item $3P_{ X } >2P_{ Y }$
  \item $P_{ X } <P_{ Y }$
\end{choices}
\end{question}
\answer{3}
